\answer{ex:06:1}{$e=(1-e^{-T_a/\tau_e})e^{-d/\tau_e}$。(a)~$0.110$对$4.5\cdot10^{-5}$，相差$\approx2.5\cdot10^3$倍。(b)~$\tau_e^*=T_a/\ln(1+T_a/d)=0.59\,\text{s}$。}
\answer{ex:06:2}{速率为每秒$\nu_x\nu_y(A_+\tau_+-A_-\tau_-)=0.8A_+$，一小时$\approx2900A_+$；$\tau_-=\tau_+/0.6=33\ms$。}
\answer{ex:06:3}{当$\mu^2+s_1^2>s_2^2$时学到$\mathbf e_1$；这里$1.01>0.25$，神经元学到$\mathbf e_1$，尽管沿$\mathbf e_2$的离散大$25$倍：规则找到的是相关矩阵而非协方差矩阵的主向量。}
\answer{ex:06:4}{灯与果汁之间$V=Re^{-(t_c+L-t)/\tau_\gamma}$，因此$\dot V=V/\tau_\gamma$；爆发面积为$Re^{-L/\tau_\gamma}$：$0.61R$和$0.78R$。}
\answer{ex:06:5}{信噪比为$1/(N+1)$，尝试次数$\propto N+1$：$5\cdot10^5$次尝试$\approx1$~个月，$5\cdot10^{11}$次尝试$\approx8\cdot10^4$~年。}
\answer{ex:06:6}{(a)~$k_2=1/\tau_d$，$k_1=1/\tau_d-1/\tau_\gamma$。(b)~虚假误差为$-(V/\tau_\gamma)\,\varepsilon/(1+\varepsilon)$，$\varepsilon=\tau_d/\tau_\gamma$，由此$\tau_d\le0.105\,\text{s}$。}
\answer{ex:06:7}{$0.46$、$0.90$、$0.99$、$0.95$，$d=3\,\text{s}$时为$0.70$。各点落在同一条关于资格迹比例$F=(1-e^{-T_{\text{cue}}/\tau_e})e^{-d/\tau_e}$的曲线上：$F\gtrsim0.1$时能学会，$F<10^{-2}$时为随机选择。}
\answer{ex:06:8}{$\eta^*=1/\bigl(2\sigma^2(N+2)\bigr)$，需$N+2$次尝试；$N$个中只有$m$个抖动时需$N(m+2)/m\ge N+2$次：稀疏化没有帮助。}
